\section{Slide 与视觉证据策略}

\subsection{Slide-first 结构化归档}
讲者在讲述 trace、agent、DSPy 优化过程时始终依赖结构化的 slide 画面来突显整体 pipeline。我们在整理过程中遵循 slide-first 原则：先捕捉静态 slide（如三类策略、APAC 宪法、DSPy 架构与 latent optimization），再在需要的地方补上动态 frame，用视觉证据锚定每一段教学逻辑。下表列出目前使用的主要视觉素材及其教学意义。
\begin{table}[H]
\centering
\begin{tabular}{lp{8cm}}
\toprule
Slide/Frame & 主要呈现内容与用途 \\
\midrule
frame-solution-categories.jpg & 三个策略（Scaling/Hybrid/Compound AI）所在的“分层” slide，帮助读者理解不同方法的成本/控制 trade-off（00:09:10--00:09:18）。 \\
frame-apac-constitution.jpg & APAC Constitution 的四个维度表格，说明判别性指标如何与 agent 提问策略绑定（00:16:08--00:16:22）。 \\
frame-dspy-architecture.jpg & DSPy 架构图，展示 latent cost、solver、plan 之间的闭环（00:32:25--00:32:45）。 \\
frame-latent-opt.jpg & 描述 latent optimization pipeline 的可微图，用来说明 Green Descent 的双层迭代（00:45:05--00:45:25）。 \\
\bottomrule
\end{tabular}
\caption{本节所用 slide/frame 与各自的教学价值}
\end{table}
\begin{knowledgebox}{Slide 与 Frame 的协同}
Slides 用于呈现结构化公式、路线图、指标表；frames 用于捕捉讲者讲解时的状态、示意图的运动或者醒目的 numerical 例子。只要 slide 可用，就优先用 slide，frame 作为补充动态或补齐解释盲点，两者一起支撑起连续的教学叙事。
\end{knowledgebox}

\subsection{关键帧时间脚注}
每张插图都严格记录其对应的时间区间，放在图注或 footer，以便读者用 subtitle 回溯细节。例如 trace curve、APAC 结构、latent optimizer 都在它们各自的小节以 \texttt{00:xx:xx--00:yy:yy} 的格式标注。若未来追加新的 slides，务必先对照 subtitles 找出最完整的显示帧，并把时间区间写清。
\begin{warningbox}{图像脚注不能省略}
缺少时间脚注会让读者无法确认 visual evidence 的来源，特别是当多个图像紧邻讲述不同主题时。每个 figure 都应该在同一页给出字幕区间，必要时在 caption 中重复一次。
\end{warningbox}

\subsection{本章小结}
Slide-first 的整理方式让抽象 pipeline 有了视觉锚点，配合 frame 的动态例证与精确时间脚注，能让读者在阅读文字时快速回到原视频的讲解现场。

